% English Version - Methodological Article
\documentclass[12pt,a4paper]{article}
\usepackage[utf8]{inputenc}
\usepackage[T1]{fontenc}
\usepackage{lmodern}
\usepackage{amsmath,amssymb}
\usepackage{graphicx}
\usepackage{xcolor}
\usepackage{hyperref}
\usepackage{geometry}
\geometry{a4paper, left=3cm, right=3cm, top=3cm, bottom=3cm}
\usepackage{setspace}
\onehalfspacing
\usepackage{parskip}
\usepackage[english]{babel}
\usepackage{csquotes}
\usepackage{microtype}
\usepackage{booktabs}
\usepackage{longtable}
\usepackage{array}
\usepackage{float}
\usepackage{url}
\usepackage{natbib}

\title{\Huge\textbf{The Empirical Grammar of Market Conversations} \\[2mm]
       \LARGE A Neuro-Symbolic Reconstruction}

\author{
  \large
  \begin{tabular}{c}
    Paul Koop
  \end{tabular}
}

\date{\large 1994--2026}

\begin{document}

\maketitle

\begin{abstract}
This paper reconstructs the empirical grammar of eight sales conversations recorded
at Aachen market square in 1994, using the Algorithmic Recursive Sequence Analysis
(ARS) framework. The analysis proceeds from raw transcripts to terminal symbol
chains, transition counting, probability induction, and finally to a probabilistic
context-free grammar (PCFG) with empirically optimized transition probabilities.
I provide an interpretive analysis of the learned probabilities, distinguishing
between \textit{constitutive rules} (structural constraints that define well-formedness)
and \textit{statistical regularities} (empirical frequencies that reflect contingent
patterns). The grammar reveals both ritualized sequences (greetings, farewells)
and strategic options (upselling loops, restarts). The paper concludes with a
discussion of how the grammar can be implemented in neuro-symbolic frameworks,
bridging qualitative interpretation and computational execution.
\end{abstract}

\newpage
\tableofcontents
\newpage

\section{Introduction: From Raw Data to Formal Grammar}

The Algorithmic Recursive Sequence Analysis (ARS) rests on a simple yet powerful
idea: social interactions leave physical traces (audio recordings) that can be
transcribed, interpreted, and transformed into formal grammars. This paper
documents the complete pipeline from raw transcripts to an empirically optimized
grammar, using eight sales conversations recorded at Aachen market square in
June/July 1994.

The contribution is threefold:

\begin{enumerate}
    \item \textbf{Empirical}: I present the full optimized grammar induced from
    the eight transcripts, with all transition probabilities.
    
    \item \textbf{Interpretive}: I analyze the learned probabilities, distinguishing
    between constitutive rules (structural necessities) and statistical regularities
    (empirical contingencies).
    
    \item \textbf{Methodological}: I discuss how the grammar can be implemented
    in neuro-symbolic frameworks, bridging qualitative interpretation and
    computational execution.
\end{enumerate}

\section{Data and Interpretation}

\subsection{The Eight Transcripts}

The empirical material comprises eight transcripts of sales conversations recorded
at Aachen market square in June/July 1994. The transcripts vary in length,
completeness, and conversation type:

\begin{table}[H]
\centering
\caption{Corpus of market conversations}
\label{tab:corpus}
\begin{tabular}{@{} l l l l @{}}
\toprule
\textbf{Text} & \textbf{Date} & \textbf{Location} & \textbf{Participants} \\
\midrule
T1 & 28.06.1994 & Butcher shop & Seller (f), Customer \\
T2 & 28.06.1994 & Cherry stall & Seller (m), C1, C2 \\
T3 & 28.06.1994 & Fish stall & Seller (m), Customer \\
T4 & 28.06.1994 & Vegetable stall & Seller (m), Customer \\
T5 & 26.06.1994 & Vegetable stall & Seller (m), C1, C2 \\
T6 & 28.06.1994 & Cheese stall & Seller (m), C1 \\
T7 & 28.06.1994 & Candy stall & Seller (m), Customer \\
T8 & 09.07.1994 & Bakery & Seller (f), Customer \\
\bottomrule
\end{tabular}
\end{table}

\subsection{Terminal Symbol Assignment}

Each utterance was assigned a terminal symbol from a predefined category system,
developed through sequential micro-analysis following objective hermeneutics
\citep{oevermann1979methodology}. The twelve terminal symbols are:

\begin{table}[H]
\centering
\caption{Terminal symbols and their meanings}
\label{tab:terminals}
\begin{tabular}{@{} l l @{}}
\toprule
\textbf{Symbol} & \textbf{Meaning} \\
\midrule
KBG & Customer greeting \\
VBG & Seller greeting \\
KBBd & Customer need (concrete) \\
VBBd & Seller inquiry \\
KBA & Customer response \\
VBA & Seller reaction \\
KAE & Customer inquiry \\
VAE & Seller information \\
KAA & Customer completion \\
VAA & Seller completion \\
KAV & Customer farewell \\
VAV & Seller farewell \\
\bottomrule
\end{tabular}
\end{table}

The resulting terminal symbol chains for the eight transcripts are:

\begin{verbatim}
T1: KBG, VBG, KBBd, VBBd, KBA, VBA, KBBd, VBBd, KBA, VAA, KAA, VAV, KAV
T2: VBG, KBBd, VBBd, VAA, KAA, VBG, KBBd, VAA, KAA
T3: KBBd, VBBd, VAA, KAA
T4: KBBd, VBBd, KBA, VBA, KBBd, VBA, KAE, VAE, KAA, VAV, KAV
T5: KAV, KBBd, VBBd, KBBd, VAA, KAV
T6: KBG, VBG, KBBd, VBBd, KAA
T7: KBBd, VBBd, KBA, VAA, KAA
T8: KBG, VBBd, KBBd, VBA, VAA, KAA, VAV, KAV
\end{verbatim}

\section{The Optimized Grammar}

\subsection{Induction Method}

The grammar was induced by counting transitions across all eight transcripts
and normalizing to probabilities:

\[
P(\sigma_j \mid \sigma_i) = \frac{\text{count}(\sigma_i \to \sigma_j)}{\sum_k \text{count}(\sigma_i \to \sigma_k)}
\]

The resulting probabilities were iteratively refined by generating artificial
chains and comparing their frequency distributions to the empirical data.
The final grammar achieved a correlation of \(r = 0.925\) between empirical
and generated frequencies.

\subsection{The Optimized Probabilistic Grammar}

\begin{table}[H]
\centering
\caption{Optimized transition probabilities}
\label{tab:grammar}
\begin{tabular}{@{} l l @{}}
\toprule
\textbf{Start symbol} & \textbf{Following symbols with probabilities} \\
\midrule
KBG & VBG (0.667), VBBd (0.333) \\
VBG & KBBd (1.0) \\
KBBd & VBBd (0.667), VAA (0.167), VBA (0.167) \\
VBBd & KBA (0.444), VAA (0.222), KBBd (0.222), KAA (0.111) \\
KBA & VBA (0.5), VAA (0.5) \\
VBA & KBBd (0.5), KAE (0.25), VAA (0.25) \\
VAA & KAA (0.857), KAV (0.143) \\
KAA & VAV (0.75), VBG (0.25) \\
VAV & KAV (1.0) \\
KAE & VAE (1.0) \\
VAE & KAA (1.0) \\
KAV & VAV (0.5), KBBd (0.5) \\
\bottomrule
\end{tabular}
\end{table}

\subsection{Graphical Representation}

Figure~\ref{fig:grammar} visualizes the grammar as a directed graph with
probabilities.

\begin{figure}[H]
\centering
\begin{verbatim}
                    KBG
                   /   \
                0.667   0.333
                 /       \
               VBG       VBBd
                |          |
               1.0         |
                |          |
               KBBd        |
               / | \        |
           0.667|  |0.167   |
             /  |  \        |
           VBBd |   VBA     |
            |   |    |      |
            |   |  0.5|     |
            |   |    |      |
            |   |   KBBd    |
            |   |   / \     |
            |   |  ... ...  |
            |   |           |
            \---/-----------/
               |
              VAA
             /   \
          0.857   0.143
           /       \
         KAA       KAV
         / \        |
      0.75 0.25     |
       /    \       |
     VAV    VBG     |
      |             |
     KAV           VAV
      |             |
      \_____________/
            |
           END
\end{verbatim}
\caption{Graphical representation of the optimized grammar. Edge weights
represent transition probabilities.}
\label{fig:grammar}
\end{figure}

\section{Interpretive Analysis of Probabilities}

\subsection{Constitutive Rules (Structural Necessities)}

Some transitions have probability 1.0, indicating that they are not merely
statistical regularities but \textit{constitutive rules} of the interaction
format:

\begin{itemize}
    \item \textbf{VBG → KBBd (1.0)}: A seller greeting must be followed by a
    customer need articulation. This is a structural property of sales
    conversations. Without this transition, the conversation would lack
    its transactional purpose.
    
    \item \textbf{KAE → VAE (1.0)}: A customer inquiry must be answered by
    seller information. This reflects the normative expectation of reciprocity
    and cooperation.
    
    \item \textbf{VAE → KAA (1.0)}: Seller information is always followed by
    customer completion. This suggests that information exchange in market
    conversations is tightly coupled to transaction closure.
    
    \item \textbf{VAV → KAV (1.0)}: Seller farewells are always reciprocated
    by customer farewells. This is a ritualized closing sequence that marks
    the end of the interaction.
\end{itemize}

\subsection{Statistical Regularities (Empirical Contingencies)}

Other transitions have probabilities between 0 and 1, reflecting empirical
frequencies that could vary across contexts:

\begin{itemize}
    \item \textbf{KBG → VBG (0.667)} vs. \textbf{KBG → VBBd (0.333)}: In
    two-thirds of cases, a customer greeting is reciprocated; in one-third,
    the seller responds directly with an inquiry (skipping the reciprocal
    greeting). This shows that the reciprocal greeting is the norm but not
    obligatory.
    
    \item \textbf{KBBd → VBBd (0.667)} vs. \textbf{KBBd → VAA (0.167)} vs.
    \textbf{KBBd → VBA (0.167)}: Most customer needs are followed by seller
    inquiries, but sometimes directly by completion (immediate purchase)
    or reaction (consultative response).
    
    \item \textbf{VBBd → KBA (0.444)} vs. \textbf{VBBd → VAA (0.222)} vs.
    \textbf{VBBd → KBBd (0.222)} vs. \textbf{VBBd → KAA (0.111)}: Seller
    inquiries have the most varied outcomes—responses, completions, need
    loops (upselling), or customer completions (early exit).
    
    \item \textbf{KBA → VBA (0.5)} vs. \textbf{KBA → VAA (0.5)}: Customer
    responses are equally likely to lead to seller reactions or direct
    completions. This suggests a strategic choice point.
    
    \item \textbf{VAA → KAA (0.857)} vs. \textbf{VAA → KAV (0.143)}: Most
    seller completions are followed by customer completions; rarely,
    directly by customer farewell (abbreviated closing).
    
    \item \textbf{KAA → VAV (0.75)} vs. \textbf{KAA → VBG (0.25)}: Three
    quarters of customer completions lead to seller farewells; one quarter
    lead to a restart of the conversation (new greeting). The restart option
    occurs when a new customer arrives or when the same customer makes
    an additional purchase.
\end{itemize}

\subsection{The Upselling Loop}

The transition \textbf{VBA → KBBd (0.5)} deserves special attention. This
loop—from seller reaction back to customer need—is the grammatical
representation of \textit{upselling}. When a seller says "Anything else?"
(VBA), the customer often responds with an additional need (KBBd).

Crucially, this loop occurs in only half of the cases (0.5). The other half
lead to completion (VAA, 0.25) or inquiry (KAE, 0.25). This suggests that
upselling is neither mandatory nor rare—it is a strategic option that sellers
can deploy, and customers can accept or deflect.

\subsection{The Restart Option}

The transition \textbf{KAA → VBG (0.25)} is particularly interesting.
A customer completion (KAA) is normally followed by farewell (VAV, 0.75),
but in a quarter of cases, it is followed by a seller greeting (VBG).
This indicates that conversations can restart after completion—for example,
when a new customer arrives (as in T2 and T5) or when the same customer
makes an additional purchase.

\subsection{The New Customer Arrival}

The transition \textbf{KAV → KBBd (0.5)} shows that a customer farewell
can be "cancelled" when a new customer arrives. In T5, a farewell (KAV)
is immediately followed by a new customer need (KBBd). This is the only
transition that breaks the otherwise strict sequencing of phases.

\section{Constitutive Rules vs. Statistical Regularities}

\subsection{The Distinction}

The ARS framework maintains a strict separation between two levels of
description:

\begin{enumerate}
    \item \textbf{Constitutive rules}: Structural constraints that define
    what counts as a well-formed sequence. These are binary (a sequence either
    conforms or it does not). In the grammar, transitions with probability
    1.0 in the corpus are candidates for constitutive rules.
    
    \item \textbf{Statistical regularities}: Empirical frequencies of
    transitions. These are probabilistic and can be updated as new data
    arrives. They reflect what \textit{happens} in the data, not what
    \textit{must} happen.
\end{enumerate}

\begin{table}[H]
\centering
\caption{Constitutive rules identified from the corpus}
\label{tab:constitutive}
\begin{tabular}{@{} l l l @{}}
\toprule
\textbf{Rule} & \textbf{Probability} & \textbf{Interpretation} \\
\midrule
VBG → KBBd & 1.0 & Seller greeting must be followed by customer need \\
KAE → VAE & 1.0 & Customer inquiry must be answered \\
VAE → KAA & 1.0 & Information must lead to completion \\
VAV → KAV & 1.0 & Farewells are reciprocal \\
\bottomrule
\end{tabular}
\end{table}

\subsection{From Statistics to Constitutivity}

A statistical regularity can become a constitutive rule through methodological
decision. For example, the transition \texttt{VBG → KBBd (1.0)} never varied
in the corpus. We could treat it as a constitutive rule: "In sales conversations,
a seller greeting must be followed by a customer need articulation." This
turns an empirical observation into a normative constraint.

However, this decision is not automatic. It requires methodological reflection:
Is this truly a necessary feature of the interaction format, or could it vary
in other contexts? The ARS approach is to treat transitions as statistical
until counterexamples force a revision of the structural grammar.

\subsection{The Methodological Significance}

This flexibility is crucial for XAI and neuro-symbolic integration. It allows
the analyst to:

\begin{enumerate}
    \item Start with purely statistical learning (no prior constraints).
    \item Identify transitions that never vary in the data.
    \item Elevate them to constitutive rules after methodological reflection.
    \item Use the resulting hybrid system for prediction and explanation.
\end{enumerate}

The system thus moves from purely empirical pattern recognition (System 1)
to rule-based explanation (System 2)—a shift that mirrors Kahneman's
distinction between fast and slow thinking \citep{kahneman2011thinking}.

\section{Toward Neuro-Symbolic Implementation}

\subsection{The Dual-Dynamics Architecture}

The ARS grammar naturally suggests a dual-dynamics architecture for
neuro-symbolic implementation:

\begin{enumerate}
    \item \textbf{Symbolic component} (System 2): The grammar rules, including
    both constitutive rules and statistical probabilities. This component is
    inspectable, falsifiable, and explainable.
    
    \item \textbf{Neural component} (System 1): A neural network that learns
    transition probabilities from data and predicts next symbols. This component
    is fast, pattern-based, and scalable.
    
    \item \textbf{Hybrid integration}: The neural component provides fast
    predictions; the symbolic component validates them against constitutive
    rules and provides explanations.
\end{enumerate}

\begin{lstlisting}[caption=Pseudocode for neuro-symbolic ARS]
# Pseudocode: ARS 5.0 Dual-Dynamics Architecture

class ARSNeuroSymbolicSystem:
    
    def __init__(self):
        # Symbolic component
        self.grammar = ARSGrammar()           # PCFG with probabilities
        self.counts = zero_matrix(12, 12)    # Transition counts
        
        # Neural component
        self.neural_network = NeuralNetwork( input_dim=12, hidden=64, output_dim=12 )
        
        # Constitutive rules (hard constraints)
        self.constitutive_rules = {
            (KBG, VBG): True, (VBG, KBBd): True,
            (KAE, VAE): True, (VAV, KAV): True, (KAV, VAV): True
        }
    
    def update_symbolic(self, from_sym, to_sym):
        """Fast, counting-based update (System 2)"""
        self.counts[from_sym][to_sym] += 1
        row_sum = sum(self.counts[from_sym])
        self.grammar.probs[from_sym] = self.counts[from_sym] / row_sum
    
    def update_neural(self, from_sym, to_sym):
        """Slow, gradient-based update (System 1)"""
        loss = cross_entropy( self.neural_network(from_sym), to_sym )
        loss.backward()
        optimizer.step()
    
    def predict_next(self, from_sym):
        """Hybrid prediction"""
        neural_probs = self.neural_network(from_sym)
        symbolic_probs = self.grammar.probs[from_sym]
        return 0.5 * neural_probs + 0.5 * symbolic_probs
    
    def is_valid(self, from_sym, to_sym):
        """Check constitutive rules"""
        return self.constitutive_rules.get((from_sym, to_sym), True)
\end{lstlisting}

\subsection{Explainability Through Proof Trees}

The symbolic component enables explainability through proof trees.
For any well-formed sequence, the grammar can produce a derivation:

\begin{lstlisting}[caption=Proof tree for a well-formed sequence]
well_formed([KBG, VBG, KBBd])
    ← start(KBG)                               [1.0]
    ← transition(KBG, VBG)                     [0.667]
    ← well_formed([VBG, KBBd])
      ← transition(VBG, KBBd)                  [1.0]
      ← well_formed([KBBd])
        ← start(KBBd)                          [0.0]
Probability: 0.667 × 1.0 × 0.0 = 0.0
\end{lstlisting}

This proof tree makes the reasoning transparent. Each step is justified by
a rule, and each rule has an explicit probability.

\subsection{Implementation Options}

The ARS grammar can be implemented in various neuro-symbolic frameworks:

\begin{table}[H]
\centering
\caption{Implementation options for ARS 5.0}
\label{tab:options}
\begin{tabular}{@{} l l l @{}}
\toprule
\textbf{Framework} & \textbf{Language} & \textbf{Best for} \\
\midrule
DeepProbLog & Prolog/Python & Research, explainability \\
PyTorch & Python & Flexibility, prototyping \\
Flux.jl & Julia & Scientific computing, performance \\
Candle & Rust & Production, edge computing \\
\bottomrule
\end{tabular}
\end{table}

Each framework has its own strengths, but all can implement the same
dual-dynamics architecture.

\section{Conclusion}

This paper has reconstructed the empirical grammar of eight market conversations,
from raw transcripts to an empirically optimized probabilistic grammar.
The analysis yielded three main insights:

\begin{enumerate}
    \item \textbf{Empirical}: The optimized grammar achieves high correlation
    with the data (\(r = 0.925\)) and reveals both constitutive rules
    (transitions with probability 1.0) and statistical regularities.
    
    \item \textbf{Interpretive}: The upselling loop (VBA → KBBd, 0.5) and
    restart option (KAA → VBG, 0.25) highlight strategic choices in sales
    interactions. The farewell-cancellation (KAV → KBBd, 0.5) shows how
    new customers can enter ongoing interactions.
    
    \item \textbf{Methodological}: The distinction between constitutive rules
    and statistical regularities provides a foundation for explainable,
    falsifiable, and adaptable neuro-symbolic systems.
\end{enumerate}

The grammar presented here is not a static artifact. It can be updated as
new data arrives (statistical plasticity) and revised when structural
anomalies are detected (structural stability). In this sense, it is a
living neuro-symbolic model—an empirical grammar that learns while
remaining explainable.

\newpage
\begin{thebibliography}{99}

\bibitem[Kahneman(2011)]{kahneman2011thinking}
Kahneman, D. (2011). \textit{Thinking, Fast and Slow}. Farrar, Straus and Giroux.

\bibitem[Koop(1994)]{koop1994scheme}
Koop, P. (1994). \textit{Grammatikinduktion empirisch gesicherter 
Verkaufsgespräche}. Scheme source code.

\bibitem[Koop(2026)]{koop2026deep}
Koop, P. (2026). \textit{From Scheme to DeepProbLog: ARS as a Methodological
Blueprint for Modern Neuro-Symbolic Programming}. the-last-freedom.org.

\bibitem[Manhaeve et al.(2018)]{manhaeve2018deepproblog}
Manhaeve, R., Dumancic, S., Kimmig, A., Demeester, T., \& De Raedt, L. (2018).
DeepProbLog: Neural probabilistic logic programming. \textit{Advances in
Neural Information Processing Systems}, 31.

\bibitem[Oevermann et al.(1979)]{oevermann1979methodology}
Oevermann, U., Allert, T., Konau, E., \& Krambeck, J. (1979). The methodology
of objective hermeneutics. In H.-G. Soeffner (Ed.), \textit{Interpretative
Procedures in the Social and Text Sciences} (pp. 352-434). Metzler.

\end{thebibliography}

\end{document}